\documentclass[12pt]{article}

\usepackage{fullpage}
\usepackage{multicol,multirow}
\usepackage{tabularx}
\usepackage{ulem}
\usepackage[utf8]{inputenc}
\usepackage{mathtools}
\usepackage{pgfplots}
\usepackage[russian]{babel}

\begin{document}

\section*{Лабораторная работа №\,9 по курсу дискрeтного анализа: Графы}

Выполнил студент группы М8О-308Б-22 МАИ \textit{Немкова Анастасия}.

\subsection*{Условие}

\textbf{Вариант 4} 

Поиск кратчайшего пути между парой вершин алгоритмом Дейкстры

Задан взвешенный неориентированный граф, состоящий из n вершин и m ребер. Вершины пронумерованы целыми числами от 1 до n. Необходимо найти длину кратчайшего пути из вершины с номером start в вершину с номером finish при помощи алгоритма Дейкстры. Длина пути равна сумме весов ребер на этом пути. Граф не содержит петель и кратных ребер.

\textbf{Формат ввода}
В первой строке заданы 1<= n <= $10^5$ и 1 <= m <= $10^5$, 1 <= start  <= n и 1 <= finish <= n. В следующих m строках записаны ребра. Каждая строка содержит три числа – номера вершин, соединенных ребром, и вес данного ребра. Вес ребра – целое число от 0 до $10^9$. 

\textbf{Формат вывода}
Необходимо вывести одно число – длину кратчайшего пути между указанными вершинами. Если пути между указанными вершинами не существует, следует вывести строку "No solution".

\subsection*{Метод решения}

Задача заключается в нахождении кратчайшего пути между двумя вершинами графа с использованием алгоритма Дейкстры. 

Для представления графа используется список смежности. Каждая вершина хранит список ребер, которые с ней соединены. Каждое ребро содержит информацию о соседней вершине и весе ребра.

Создаём массив расстояний, где для каждой вершины будет храниться минимальное расстояние от исходной вершины. Изначально в массиве значения расстояний до каждой вершины устанавливаются как бесконечность, за исключением расстояния до стартовой вершины, которое равно нулю.

Создаём приоритетную очередь, в которую будут добавляться непосещенные вершины. 

На каждом шаге извлекаем вершину с минимальным расстоянием из приоритетной очереди и для каждой смежной с ней вершины производим релаксацию ребра, т. е. проверяем, можно ли улучшить её расстояние через текущую вершину (сумма расстояния до текущей вершины и веса рёбра, ведущего в смежную вершину, должна быть меньше текущего расстояния до смежной вершины). Обновляем расстояние до смежной вершины и добавляем её в приоритетную очередь. 

Таким образом обрабатываем все вершины и получаем кратчайшие расстояния от исходной вершины до всех остальных вершин.

\subsection*{Описание программы}

Программа реализует алгоритм Дейкстры для нахождения кратчайшего пути между двумя вершинами в графе с неотрицательными весами рёбер. 

Вводятся количество вершин n, количество рёбер m, начальная вершина start и конечная вершина finish. После этого программа строит граф, используя список смежности, и применяет алгоритм Дейкстры для нахождения минимальных расстояний от начальной вершины до всех остальных вершин. В конце выводится кратчайшее расстояние от вершины start до вершины finish, если путь существует, иначе выводится сообщение "No solution".

\textbf{Пространственная сложность}

Граф представлен списком смежности, для которого требуется O(n + m) памяти, где n — количество вершин, а m — количество рёбер.

Массив для хранения расстояний от начальной вершины к остальным вершинам требует O(n) памяти.

Множество, используемое для выбора минимальной вершины, хранит в себе максимум O(n) элементов.

Таким образом, общая пространственная сложность программы составляет O(n + m).

\textbf{Временная сложность}

Для каждой вершины выполняются операции вставки и удаления в множестве с приоритетом, что занимает время O(log n).

Таким образом, общая временная сложность алгоритма Дейкстры составляет O((n + m) * log n), где n — количество вершин, а m — количество рёбер в графе.

\subsection*{Дневник отладки}

\begin{enumerate}
    \item 18 дек 2024, 16:57:44 OK 

    Использование set, вместо очереди с приоритетом


    \item 18 дек 2024, 18:00:54 WA на 8 тесте

    Была произведена попытка сделать алгоритм для неразряженного графа

\end{enumerate}


\subsection*{Тест производительности}

Для измерения производительсти сравним время работы алгоритма Дейкстры и поиска в глубину(DFS), в котором ,начиная с заданной вершины, рекурсивно проходим по графу, обновляя минимальное расстояние до каждой вершины, если найден путь короче текущего. 

Тесты состоят из плотных и разряженных графов с $10^4$ вершин и от $10^$ до $10^5$ ребер

\begin{figure}[htbp]
    \centering
    \begin{tikzpicture}
        \begin{semilogyaxis}[
            xlabel={Количество рёбер ($m \times 10^3$)},
            ylabel={Время работы (секунды)},
            grid=major,
            xmin=1, xmax=10,
            ymin=0.00001, ymax=10,
            xtick={1, 2, 3, 4, 5, 6, 7, 8, 9, 10},
            ytick={0.00001, 0.0001, 0.001, 0.01, 0.1, 1, 10},
            legend style={at={(0.5,-0.2)},anchor=north},
            legend columns=2,
            width=0.8\textwidth,
            height=0.5\textwidth,
        ]
        \addplot[color=blue,mark=*] coordinates {
            (1, 0.000320998) 
            (2, 0.000802995)
            (3, 0.00132936)
            (4, 0.00171742)
            (5, 0.00190863)
            (6, 0.00207361)
            (7, 0.00252237)
            (8, 0.00292437)
            (9, 0.0036536)
            (10, 0.00406567)
        };
        \addlegendentry{Алгоритм Дейкстры}
        
        \addplot[color=red,mark=square] coordinates {
            (1, 9.5715e-05)
            (2, 0.0296836)
            (3, 0.103817)
            (4, 0.175599)
            (5, 0.23708)
            (6, 0.342011)
            (7, 0.449766)
            (8, 0.565417)
            (9, 0.694565)
            (10, 0.780594)
        };
        \addlegendentry{DFS (Наивный)}
        \end{semilogyaxis}
    \end{tikzpicture}
    \label{fig:graph}
\end{figure}

Из графика видим, что алгоритм Дейкстры имеет логарифмическую сложность O((n + m) * log n), а DFS - линейную O(n + m), где n — количество вершин, а m — количество рёбер в графе.


\subsection*{Выводы}

В ходе данной лабораторной работы был изучен алгоритм Дейкстры, предназначенный для нахождения кратчайшего пути в графе с положительными весами рёбер. 

Алгоритм эффективно работает с графами, используя такую структуру данных, как приоритетная очередь, что обеспечивает логарифмическую сложность операций. Также было проведено сравниение алгоритма Дейкстры и метода поиска в глубину (DFS), где алгоритм Дейкстры показал значительно лучшие результаты при обработке графов с большим количеством рёбер. Алгоритм Дейкстры находит широкое применение в реальных задачах, таких как навигационные системы и сети связи, где важно находить оптимальные маршруты.

\end{document}





